\documentclass[UTF8,fontset=fandol,zihao=-4]{ctexart}

\usepackage[a4paper,top=22mm,bottom=22mm,left=21mm,right=21mm,headheight=24pt]{geometry}
\usepackage[table]{xcolor}
\usepackage{booktabs}
\usepackage{tabularx}
\usepackage{longtable}
\usepackage{array}
\usepackage{ragged2e}
\usepackage{enumitem}
\usepackage[most]{tcolorbox}
\usepackage{fancyhdr}
\usepackage{lastpage}
\usepackage{xurl}
\usepackage[backend=biber,style=numeric,sorting=nyt,giveninits=true]{biblatex}
\addbibresource{references.bib}
\renewcommand*{\mkbibacro}[1]{\MakeUppercase{#1}}
\AtBeginBibliography{\pagestyle{fancy}\thispagestyle{fancy}}
\usepackage{hyperref}
\usepackage{bookmark}

\definecolor{navy}{HTML}{17324D}
\definecolor{blue}{HTML}{2B6CB0}
\definecolor{teal}{HTML}{0F766E}
\definecolor{orange}{HTML}{C66A13}
\definecolor{redbrown}{HTML}{9B3A34}
\definecolor{ink}{HTML}{243447}
\definecolor{muted}{HTML}{64748B}
\definecolor{line}{HTML}{D8E1EA}
\definecolor{paperblue}{HTML}{F3F7FB}
\definecolor{paperorange}{HTML}{FFF7ED}

\newcommand{\ReportTitle}{__REPORT_TITLE__}
\newcommand{\ReportSubtitle}{__REPORT_SUBTITLE__}
\newcommand{\ReportAuthor}{__REPORT_AUTHOR__}
\newcommand{\ReportDate}{__REPORT_DATE__}
\newcommand{\ReportScope}{__REPORT_SCOPE__}

\hypersetup{
  colorlinks=true,
  linkcolor=navy,
  urlcolor=blue,
  citecolor=teal,
  pdftitle={\ReportTitle},
  pdfauthor={\ReportAuthor}
}

\ctexset{
  section={format=\Large\bfseries\sffamily\color{navy},beforeskip=1.5ex,afterskip=0.8ex},
  subsection={format=\large\bfseries\sffamily\color{blue},beforeskip=1.2ex,afterskip=0.5ex},
  subsubsection={format=\normalsize\bfseries\sffamily\color{teal}}
}

\pagestyle{fancy}
\fancyhf{}
\fancyhead[L]{\sffamily\small\color{muted}\ReportTitle}
\fancyhead[R]{\sffamily\small\color{muted}\ReportDate}
\fancyfoot[C]{\sffamily\small\color{muted}第 \thepage\ 页，共 \pageref{LastPage}\ 页}
\renewcommand{\headrulewidth}{0.4pt}
\fancypagestyle{plain}{
  \fancyhf{}
  \fancyhead[L]{\sffamily\small\color{muted}\ReportTitle}
  \fancyhead[R]{\sffamily\small\color{muted}\ReportDate}
  \fancyfoot[C]{\sffamily\small\color{muted}第 \thepage\ 页，共 \pageref{LastPage}\ 页}
  \renewcommand{\headrulewidth}{0.4pt}
}
% ctex may finalize its class page styles at begin-document time. Re-apply the
% report's plain style afterwards so longtable and bibliography continuation
% pages do not fall back to the class-default footer.
\AtBeginDocument{
  \fancypagestyle{plain}{
    \fancyhf{}
    \fancyhead[L]{\sffamily\small\color{muted}\ReportTitle}
    \fancyhead[R]{\sffamily\small\color{muted}\ReportDate}
    \fancyfoot[C]{\sffamily\small\color{muted}第 \thepage\ 页，共 \pageref{LastPage}\ 页}
    \renewcommand{\headrulewidth}{0.4pt}
  }
}

\setlength{\parindent}{0pt}
\setlength{\parskip}{0.45em}
\setlength{\tabcolsep}{5pt}
\renewcommand{\arraystretch}{1.22}
\setlist[itemize]{leftmargin=1.6em,itemsep=0.25em,topsep=0.3em}
\setlist[enumerate]{leftmargin=1.8em,itemsep=0.25em,topsep=0.3em}

\newcolumntype{L}[1]{>{\RaggedRight\arraybackslash}p{#1}}
\newcolumntype{Y}{>{\RaggedRight\arraybackslash}X}

\newcommand{\claimref}[1]{\textsuperscript{\colorbox{blue!10}{\textcolor{blue}{\scriptsize #1}}}}
\newcommand{\ReportPlaceholder}[1]{\textcolor{redbrown}{[请替换：#1]}}

\newtcolorbox{summarybox}{
  enhanced,breakable,colback=paperblue,colframe=blue!35,
  boxrule=0.6pt,arc=2mm,left=3mm,right=3mm,top=2mm,bottom=2mm
}

\newtcolorbox{scopebox}{
  enhanced,breakable,colback=paperorange,colframe=orange!45,
  boxrule=0.6pt,arc=2mm,left=3mm,right=3mm,top=2mm,bottom=2mm
}

\begin{document}
\hypersetup{pageanchor=false}

\begin{titlepage}
\thispagestyle{empty}
\vspace*{5mm}
\begin{tcolorbox}[
  enhanced,colback=navy,colframe=navy,boxrule=0pt,arc=1.5mm,
  left=9mm,right=9mm,top=11mm,bottom=11mm,
  borderline south={2.5pt}{0pt}{blue}
]
{\color{white}\sffamily\bfseries\fontsize{23}{31}\selectfont
\ReportTitle\par}
\vspace{3mm}
{\color{white}\sffamily\fontsize{14}{21}\selectfont
\ReportSubtitle\par}
\end{tcolorbox}

\vspace{16mm}
\begin{summarybox}
{\sffamily\bfseries\color{navy}一句话结论}\par
\ReportPlaceholder{用证据边界清晰的一句话总结当前能支持的判断。}
\end{summarybox}

\vspace{8mm}
\begin{tabularx}{\textwidth}{L{30mm}Y}
\textcolor{muted}{研究范围} & \ReportScope\\
\textcolor{muted}{整理者} & \ReportAuthor\\
\textcolor{muted}{访问与整理日期} & \ReportDate\\
\textcolor{muted}{证据覆盖} & \ReportPlaceholder{正文已读、部分读取、仅元信息与不可读来源数量。}\\
\end{tabularx}

\vfill
{\sffamily\small\color{muted}
本报告将来源明确说明、实际观察、综合推断、待验证假设和证据不足分开表达。}
\end{titlepage}

\hypersetup{pageanchor=true}
\tableofcontents
\clearpage

\section{执行摘要}

\begin{summarybox}
\begin{enumerate}
  \item \ReportPlaceholder{关键判断一，并在句末添加 Claim ID。}
  \item \ReportPlaceholder{关键判断二，并保留适用范围。}
  \item \ReportPlaceholder{最重要的证据缺口和下一项决定性检查。}
\end{enumerate}
\end{summarybox}

\section{研究问题、来源与方法}

\subsection{研究问题与边界}

\ReportPlaceholder{说明问题、读者、时间或版本、明确不回答的内容。}

\subsection{来源准入与覆盖}

\begin{scopebox}
\ReportPlaceholder{说明主来源、补充来源、读取深度、不可读来源和未独立复现的边界。}
\end{scopebox}

\section{证据矩阵}

\begin{longtable}{L{16mm}L{30mm}L{30mm}L{25mm}L{35mm}}
\caption{关键 Claim 与证据}\label{tab:claim-evidence}\\
\toprule
Claim & 判断 & 类型 & 来源 & 边界\\
\midrule
\endfirsthead
\multicolumn{5}{l}{\small\color{muted}表 \thetable{}（续）}\\
\toprule
Claim & 判断 & 类型 & 来源 & 边界\\
\midrule
\endhead
\midrule
\multicolumn{5}{r}{\small\color{muted}下页续}\\
\endfoot
\bottomrule
\endlastfoot
C001 & \ReportPlaceholder{完整 Claim} & \ReportPlaceholder{source statement 等} & S001 & \ReportPlaceholder{适用范围}\\
\end{longtable}

\section{主要发现}

\subsection{来源明确说明或实际观察}

\ReportPlaceholder{先写直接证据。关键句后使用 claimref，例如：结论。claimref\{C001\}。}

\subsection{综合推断}

\ReportPlaceholder{说明推导链、替代解释和适用范围。}

\section{冲突、不可比项与证据缺口}

\ReportPlaceholder{并列冲突来源、版本和口径；列出最小补证动作。}

\section{结论与下一项决定性检查}

\ReportPlaceholder{区分当前结论、待验证假设、下一项检查和停止条件。}

\appendix
\section{来源台账}

\begin{longtable}{L{14mm}L{38mm}L{34mm}L{23mm}L{34mm}}
\caption{来源、读取状态与证据范围}\label{tab:source-ledger}\\
\toprule
ID & 标题 & 版本/定位 & 读取深度 & 可支持范围与限制\\
\midrule
\endfirsthead
\multicolumn{5}{l}{\small\color{muted}表 \thetable{}（续）}\\
\toprule
ID & 标题 & 版本/定位 & 读取深度 & 可支持范围与限制\\
\midrule
\endhead
\bottomrule
\endlastfoot
S001 & \ReportPlaceholder{来源标题} & \ReportPlaceholder{版本与定位} & full/partial & \ReportPlaceholder{范围与限制}\\
\end{longtable}

\printbibliography[heading=bibintoc,title={参考文献}]

\end{document}
